\documentclass[../../main.tex]{subfiles}

\begin{document}

\chapter[Pengenalan Kompilasi dan Arsitektur Kompilator]{Pengenalan Kompilasi dan Arsitektur Kompilator}
\label{chap:arsitektur}

\begin{outcomebox}
\textbf{Sub-CPMK:} \subcpmk{1.1}, \subcpmk{1.2}, \subcpmk{1.3}\\[3pt]
Setelah menyelesaikan bab ini, mahasiswa mampu:
\begin{enumerate}[label=\arabic*.,leftmargin=*]
  \item membedakan translator, assembler, interpreter, dan compiler;
  \item menjelaskan fase-fase kompilasi beserta peran tabel simbol dan
        penanganan kesalahan;
  \item menganalisis trade-off kompilator satu lintasan dan banyak
        lintasan;
  \item menjelaskan pengaruh lingkungan run-time terhadap organisasi
        memori program;
  \item menyiapkan struktur proyek mini-kompilator Subset-C (M0).
\end{enumerate}
\end{outcomebox}

\section*{Materi Pokok}
\begin{itemize}[leftmargin=*]
  \item Keluarga translator: assembler, compiler, interpreter, dan
        cross-compiler.
  \item Fase analisis dan sintesis pada kompilator.
  \item Manajemen tabel simbol dan penanganan kesalahan.
  \item Kompilator satu lintasan versus banyak lintasan.
  \item Lingkungan run-time: alokasi statis, stack, dan heap.
\end{itemize}

\subfile{section-01}
\subfile{section-02}
\subfile{section-03}
\subfile{section-04}
\subfile{section-05}

\section*{Contoh}
\addcontentsline{toc}{section}{Contoh}

\begin{contoh}[Aliran token program Subset-C]
Diberikan program
\begin{lstlisting}[language=C]
int main() {
    int a = 5;
    print(a);
    return 0;
}
\end{lstlisting}
Aliran token yang dihasilkan secara konseptual adalah:
\begin{center}
\small
\code{INT} \code{ID(main)} \code{LPAREN} \code{RPAREN} \code{LBRACE}
\code{INT} \code{ID(a)} \code{ASSIGN} \code{INT\_LIT(5)} \code{SEMI}
\code{PRINT} \code{LPAREN} \code{ID(a)} \code{RPAREN} \code{SEMI}
\code{RETURN} \code{INT\_LIT(0)} \code{SEMI} \code{RBRACE} \code{EOF}.
\end{center}
Setiap token membawa lokasi (baris, kolom) untuk keperluan diagnostik.
\end{contoh}

\begin{contoh}[Compiler versus interpreter]
Program yang sama dijalankan pada compiler dan interpreter. Pada
compiler, program sumber dikompilasi menjadi executable, lalu
executable dijalankan berulang kali. Pada interpreter, program sumber
dibaca dan dieksekusi setiap kali dijalankan. Jika program dijalankan
100 kali, compiler membayar biaya kompilasi sekali, sedangkan
interpreter membayar biaya parsing dan evaluasi 100 kali. Namun, untuk
program yang hanya dijalankan sekali, interpreter sering lebih cepat
karena tidak perlu menghasilkan kode target.
\end{contoh}

\begin{contoh}[Activation record]
Misalkan fungsi \code{faktorial(3)} dipanggil. Frame pertama memuat
\code{n = 3} dan alamat kembali ke pemanggil. Karena \code{n > 1},
fungsi memanggil \code{faktorial(2)}, mendorong frame kedua dengan
\code{n = 2}. Proses berlanjut hingga \code{faktorial(1)}. Setelah
basis tercapai, nilai 1 dikembalikan, lalu frame dilepas satu per
satu sambil mengalikan hasilnya: \(1 \times 2 \times 3 = 6\).
\end{contoh}

\begin{contoh}[Memilih strategi kompilasi]
Sebuah tim mengembangkan bahasa domain khusus untuk perangkat
mikrokontroler dengan memori sangat terbatas. Mereka memilih
kompilator satu lintasan yang langsung menghasilkan kode mesin.
Keputusan ini masuk akal karena memori kecil dan optimasi agresif
tidak menjadi prioritas. Sebaliknya, tim yang mengembangkan bahasa
umum untuk banyak platform memilih front-end bersama dan beberapa
back-end, sehingga investasi pada analisis dan IR terbayar.
\end{contoh}

\section*{Aktivitas Pembelajaran dan Praktikum}
\addcontentsline{toc}{section}{Aktivitas Pembelajaran dan Praktikum}

\begin{table}[H]
\centering
\caption{Rencana aktivitas Bab III}
\label{tab:ch03-aktivitas}
\begin{tabularx}{\textwidth}{@{}llX@{}}
\toprule
\textbf{Sesi} & \textbf{Bentuk} & \textbf{Kegiatan} \\
\midrule
1 & Ceramah & Keluarga translator dan perbedaan compiler-interpreter \\
2 & Diskusi & Fase kompilasi dan aliran informasi \\
3 & Latihan & Membaca tabel simbol dan pesan kesalahan \\
4 & Praktikum & M0: menyiapkan repositori dan spesifikasi Subset-C \\
\bottomrule
\end{tabularx}
\end{table}

\subsection*{M0: Repositori dan Spesifikasi Bahasa}
\textbf{Tujuan.} Menyiapkan struktur proyek dan spesifikasi awal
Subset-C sebagai fondasi seluruh milestone berikutnya.

\textbf{Langkah kerja.}
\begin{enumerate}[label=\arabic*.,leftmargin=*]
  \item Salin direktori \file{code/subsetc} ke repositori pribadi.
  \item Baca \file{code/subsetc/spec/subset-c.md} dan
        \file{code/subsetc/README.md}.
  \item Jalankan \file{build.bat} atau \file{Makefile} untuk memeriksa
        ketersediaan toolchain. Jika toolchain belum ada, catat
        langkah instalasinya.
  \item Tambahkan minimal dua program contoh baru ke
        \file{examples/}: satu program aritmetika dan satu program
        dengan kontrol alur.
  \item Tuliskan satu halaman catatan tentang keputusan desain bahasa
        yang Anda ambil.
\end{enumerate}

\textbf{Definisi selesai.}
\begin{checklist}
  \item Repositori memiliki struktur \file{spec}, \file{include},
        \file{src}, \file{examples}, \file{tests}.
  \item Spesifikasi v0.1 terbaca dan dipahami.
  \item Minimal tiga program contoh tersedia.
  \item Catatan keputusan desain diserahkan.
\end{checklist}

\textbf{Penilaian.} M0 tidak memiliki bobot tersendiri, tetapi menjadi
prasyarat Praktikum 1 (8\%). Rubrik yang dipakai adalah rubrik umum
praktikum pada Lampiran A.

\section*{Latihan dan Refleksi}
\addcontentsline{toc}{section}{Latihan dan Refleksi}

\subsection*{Latihan}
\begin{enumerate}[label=\arabic*.,leftmargin=*]
  \item Jelaskan perbedaan assembler, compiler, interpreter, dan
        cross-compiler. Berikan satu contoh untuk masing-masing.
  \item Sebutkan enam fase kompilasi dan jelaskan masukan serta
        keluaran setiap fase.
  \item Mengapa tabel simbol disebut pusat informasi kompilator?
  \item Berikan tiga contoh kesalahan leksikal, tiga kesalahan
        sintaksis, dan tiga kesalahan semantik.
  \item Diberikan program C sederhana, tuliskan aliran token beserta
        lokasi setiap token.
  \item Gambarkan pohon sintaks untuk program tersebut.
  \item Tuliskan kode antara tiga alamat untuk program tersebut.
  \item Lengkapi tabel penelusuran fase dan jelaskan satu sel yang
        paling sulit Anda isi.
\end{enumerate}

\subsection*{Refleksi}
\begin{reflectionbox}
\begin{enumerate}[label=\arabic*.,leftmargin=*]
  \item Setelah membaca bab ini, fase mana yang paling menarik bagi
        Anda? Mengapa?
  \item Apakah Anda dapat menjelaskan mengapa compiler dan interpreter
        dapat menghasilkan perilaku yang berbeda untuk program yang
        sama?
  \item Apa satu hal yang akan Anda lakukan berbeda saat membaca
        program setelah memahami token, pohon sintaks, dan kode antara?
\end{enumerate}
\end{reflectionbox}

\section*{Asesmen}
\addcontentsline{toc}{section}{Asesmen}

Asesmen Bab III mencakup Kuis (5\%) dan bagian awal Tugas Teori
(8\%). Fokusnya adalah pemahaman arsitektur dan aliran informasi
antar-fase.

\begin{table}[H]
\centering
\caption{Instrumen asesmen Bab III}
\label{tab:ch03-asesmen}
\begin{tabularx}{\textwidth}{@{}lXl@{}}
\toprule
\textbf{Instrumen} & \textbf{Indikator} & \textbf{Bobot} \\
\midrule
Kuis & Membedakan translator, compiler, interpreter, dan assembler &
5\% \\
Kuis & Menjelaskan fase kompilasi dan peran tabel simbol & 5\% \\
Tugas Teori & Menganalisis trade-off one-pass dan multi-pass & 8\% \\
Praktikum M0 & Repositori dan spesifikasi v0.1 lengkap & prasyarat \\
\bottomrule
\end{tabularx}
\end{table}

Kriteria kelulusan: mahasiswa dapat menjelaskan pipeline kompilasi,
menempatkan tabel simbol dan penanganan kesalahan pada pipeline, serta
memilih strategi kompilasi yang sesuai dengan kebutuhan.

\section*{Checklist Pencapaian Kompetensi}
\addcontentsline{toc}{section}{Checklist Pencapaian Kompetensi}

\begin{checklist}
  \item Saya dapat membedakan translator, assembler, compiler, dan
        interpreter.
  \item Saya dapat menjelaskan perbedaan compiler dan interpreter.
  \item Saya dapat menyebutkan fase analisis dan sintesis.
  \item Saya dapat menjelaskan peran tabel simbol pada setiap fase.
  \item Saya dapat mengklasifikasikan kesalahan program.
  \item Saya dapat menulis pesan kesalahan yang informatif.
  \item Saya dapat menjelaskan organisasi front-end dan back-end.
  \item Saya dapat menyusun aliran token untuk program C sederhana.
  \item Saya dapat menggambarkan pohon sintaks untuk program C
        sederhana.
  \item Saya dapat menuliskan kode antara tiga alamat.
  \item Saya dapat mengisi tabel penelusuran fase: token, pohon
        sintaks, dan kode antara.
  \item Saya sudah menyelesaikan M0: repositori dan spesifikasi v0.1.
\end{checklist}

\section*{Rangkuman}
\addcontentsline{toc}{section}{Rangkuman}

Pada bab ini telah dibahas:
\begin{itemize}[leftmargin=*]
  \item keluarga translator dan perbedaan compiler, interpreter,
        assembler, serta cross-compiler;
  \item fase analisis (leksikal, sintaksis, semantik) dan sintesis
        (IR, optimasi, pembangkitan kode);
  \item peran tabel simbol dan strategi penanganan kesalahan;
  \item penelusuran fase: token, pohon sintaks, dan kode antara;
  \item cara membaca dan memeriksa keluaran antara setiap fase.
\end{itemize}

Bab IV memperdalam landasan formal yang dipakai pada setiap fase,
khususnya hierarki Chomsky dan hubungan antara kelas bahasa dengan
mesin pengenal. Pemahaman ini akan menjelaskan mengapa lexer cukup
memakai otomata hingga, sedangkan parser memerlukan otomata pushdown.

\section*{Evaluasi}
\addcontentsline{toc}{section}{Evaluasi}

\begin{enumerate}[label=\arabic*.,leftmargin=*]
  \item Jelaskan perbedaan mendasar antara compiler dan interpreter,
        lalu berikan satu situasi ketika interpreter lebih tepat
        dipilih.
  \item Gambarkan pipeline kompilasi dan jelaskan perubahan
        representasi pada setiap fase.
  \item Mengapa analisis semantik memerlukan tabel simbol, sedangkan
        analisis leksikal tidak?
  \item Diberikan program dengan kesalahan pada baris 5, 8, dan 12,
        rancang strategi pelaporan agar pengguna dapat memperbaiki
        ketiganya sekaligus.
  \item Bandingkan one-pass dan multi-pass untuk bahasa yang
        mendukung fungsi rekursif dan pemanggilan sebelum deklarasi.
  \item Jelaskan bagaimana pemisahan front-end dan back-end membantu
        portabilitas compiler.
  \item Gambarkan activation record untuk fungsi
        \code{int f(int a, int b) \{ int x; int y; return a + b + x + y; \}}.
\end{enumerate}

\subsection*{Kunci dan Rubrik Singkat}
\begin{itemize}[leftmargin=*]
  \item Soal 1--3: ketepatan konsep (60\%), contoh (40\%).
  \item Soal 4: kelengkapan strategi (50\%), kualitas pesan (50\%).
  \item Soal 5--6: analisis trade-off (70\%), kesimpulan (30\%).
  \item Soal 7: kebenaran struktur frame (80\%), kelengkapan (20\%).
\end{itemize}


\ifSubfilesClassLoaded{%
  \bibliographystyle{apalike}%
  \bibliography{../../references}%
}{}

\end{document}
